\section{From boundary values to fields on a common domain}
\label{sec:word-reconstruction}
The Euclidean Hilbert space and the relativistic field Hilbert space have
related, but different, definitions. The bridge between them is a limit of
vectors with small positive imaginary times. This section proves the bridge and
the ensuing operator construction. The preceding two sections construct the analytic boundary values and
verify their heat-vector bridge from the stated slab hypotheses. Positivity
of the boundary values will now follow from that bridge.

Write $E$ for the algebraic union of the finite-dimensional field spaces and
$\mathscr S=\mathscr S(\mathbb R^4;\mathbb C)$ for the Schwartz space. A letter is
a pair $a=(v,f)\in E\times\mathscr S$, and a word is a finite ordered list of
letters. Let $J:E\to E$ be an antilinear involution, put
$a^*=(Jv,\overline f)$, and define the star of a word by reversing its order
and starring each letter. The empty word is denoted by $1$. Concatenation is
written without an intervening symbol.

Suppose $W_n(v_1,\ldots,v_n)$ is a tempered distribution on
$(\mathbb R^4)^n$, multilinear in the labels. Its value on a product test
function defines $w(a_1\cdots a_n)$. Set
\[
 K(b,c)=w(b^*c).
\]
The joint distribution, rather than merely this restricted family of tests,
will be needed in constructing the boundary values. For the operator argument
below the following properties specify exactly what is used.
\begin{assumption}[Boundary data]\label{input:word-boundary}
The functional $w$ is multilinear separately in every label and every test
function, and $w(1)=1$. For fixed words $b,c$ and labels $v,u$, the function
\[
 (f,g)\longmapsto K((v,f)b,(u,g)c)
\]
is jointly continuous in the Schwartz topology. There are vectors
$h_t(b)$, $t>0$, in the selected Euclidean Hilbert space $\mathcal H_E$,
with unit vacuum $\Omega_E$, such that
\begin{align}
 \lim_{s,t\downarrow0}\langle h_s(b),h_t(c)\rangle&=K(b,c),
       \label{eq:independent-heat-gram}\\
 \lim_{t\downarrow0}h_t(1)&=\Omega_E.\nonumber
\end{align}
The two positive parameters in the first limit vary independently.
\end{assumption}
Here and below the inner product is antilinear in its first argument.
In the application $h_t(b)$ is the integral of a product of bounded slab
insertions and intervening complex-time translations, at strictly increasing
heights $(t,2t,\ldots,nt)$. Thus the integrability of this vector-valued
integral and the joint limit \eqref{eq:independent-heat-gram} require proof.
A scalar boundary limit along one fixed ray does not by itself give this input.

\begin{lemma}[Removing the imaginary heights]\label{lem:cold-vectors}
Under Input~\ref{input:word-boundary}, each $h_t(b)$ has a norm limit
$q(b)\in\mathcal H_E$, and
\[
 \langle q(b),q(c)\rangle=K(b,c),\qquad q(1)=\Omega_E.
\]
In particular $K$ is a positive semidefinite kernel.
\end{lemma}
\begin{proof}
Fix $b$. Expanding a squared norm gives
\[
 \|h_s(b)-h_t(b)\|^2
 =\langle h_s(b),h_s(b)\rangle+\langle h_t(b),h_t(b)\rangle
  -\langle h_s(b),h_t(b)\rangle-\langle h_t(b),h_s(b)\rangle.
\]
All four terms tend to $K(b,b)$, including the diagonal terms, by the joint
limit. The expression therefore tends to zero as $s,t\downarrow0$.
Completeness gives the norm limit. Continuity of the inner product gives the
asserted Gram identities and the specified empty-word limit gives the vacuum.
For any finitely many coefficients $c_i$ and words $b_i$,
\[
 \sum_{i,j}\overline{c_i}K(b_i,b_j)c_j
   =\Big\|\sum_i c_iq(b_i)\Big\|^2\geq0.
\]
\end{proof}

Let $\mathcal P$ be the vector space of formal finite complex linear
combinations of words. Extend $K$ sesquilinearly to $\mathcal P$ and let
$\mathcal N=\{p:K(p,p)=0\}$. Cauchy--Schwarz for a positive semidefinite form
follows by applying positivity to $p+zq$ and minimizing over $z$ (the case
$K(q,q)=0$ follows by allowing $|z|$ to grow). Consequently
$\mathcal N$ is the radical of the form. The quotient
$\mathcal D=\mathcal P/\mathcal N$ is an inner-product space; let
$\mathcal H$ be its completion and $[b]\in\mathcal D$ the class of a word.

\begin{proposition}[The common invariant domain]\label{prop:word-fields}
There is a bilinear map
\[
 (v,f)\longmapsto\Phi(v,f):\mathcal D\to\mathcal D,
 \qquad \Phi(v,f)[b]=[(v,f)b].
\]
The vacuum $\Omega=[1]$ has norm one, the domain $\mathcal D$ is dense and
spanned by its word vectors, and
\begin{align}
 \langle\Phi(v,f)\xi,\eta\rangle
   &=\langle\xi,\Phi(Jv,\overline f)\eta\rangle,
                       \label{eq:word-adjoint}\\
 \langle\Omega,\Phi(a_1)\cdots\Phi(a_n)\Omega\rangle
   &=w(a_1\cdots a_n).\nonumber
\end{align}
For each $v$ and $\xi\in\mathcal D$, the map
$f\mapsto\Phi(v,f)\xi$ is continuous from $\mathscr S$ to $\mathcal H$.
As densely defined operators on $\mathcal H$,
$\Phi(Jv,\overline f)\subseteq\Phi(v,f)^*$; in particular every field
operator is closable.
\end{proposition}
\begin{proof}
Let $L_a$ denote left concatenation on formal sums. Word reversal gives the
exact identity
\[
 K(L_ap,q)=K(p,L_{a^*}q).
\]
If $p\in\mathcal N$, the right side vanishes for every $q$. Taking
$q=L_ap$ proves $K(L_ap,L_ap)=0$. Thus $L_a$ preserves $\mathcal N$ and
descends to $\mathcal D$. This is the necessary null-space argument;
positivity without the concatenation identity would not suffice.
The same identity proves \eqref{eq:word-adjoint}. Multilinearity of $w$
shows that every matrix element against a word of the defect in a proposed
linearity identity is zero. Since words span $\mathcal D$, each such defect
is zero. This proves bilinearity of $\Phi$ even though no multilinearity
relations were imposed on the formal word space initially.

The vacuum norm, cyclicity and moment identity follow directly from the
construction. To prove continuity first take $\xi=[b]$ and write
$\eta_f=[(v,f)b]$. At any $f_0$,
\[
 \|\eta_f-\eta_{f_0}\|^2
 =K((v,f)b,(v,f)b)+K((v,f_0)b,(v,f_0)b)
  -2\operatorname{Re}K((v,f_0)b,(v,f)b),
\]
which tends to zero by the assumed joint continuity. Finite linear
combinations give the result for every $\xi\in\mathcal D$.
Finally \eqref{eq:word-adjoint}, with $\mathcal D$ dense, is precisely the
criterion for the asserted adjoint inclusion. The adjoint domain contains the
dense space $\mathcal D$; hence the original operator is closable. Equivalently,
if $\xi_n\to0$ and $\Phi(v,f)\xi_n\to\eta$, pairing with each vector of
$\mathcal D$ gives $\eta=0$, which directly proves closability.
\end{proof}

\begin{proposition}[Locality and retention of labels]\label{prop:word-locality}
Suppose, whenever $f,g$ have mutually spacelike supports, that
\[
 w(b(v,f)(u,g)c)=w(b(u,g)(v,f)c)
\]
for all exterior words $b,c$. Then
$[\Phi(v,f),\Phi(u,g)]=0$ on $\mathcal D$.
Suppose also that a label $v$ whose insertion has zero expectation against
all exterior words and all tests must be zero. Then for every $v\ne0$
there is an $f$ with $\Phi(v,f)\ne0$.
\end{proposition}
\begin{proof}
Pair the commutator applied to $[c]$ with $[b]$. Its value is the difference
of the two assumed expectations, with left exterior word $b^*$, and so is
zero. Density and linearity prove the operator identity on $\mathcal D$.
If all $\Phi(v,f)$ vanish, then for any $b,c,f$,
\[
 w(b(v,f)c)=\langle[b^*],\Phi(v,f)[c]\rangle=0,
\]
because $(b^*)^*=b$. The detection assumption now gives $v=0$.
\end{proof}

\subsection{The group action and its Euclidean embedding}
Let $G$ be the Poincar\'e group with proper orthochronous Lorentz part. Suppose its action on labels
and pullback action on test functions induce an action $b\mapsto gb$ on
words, preserving the empty word. Assume
\begin{equation}\label{eq:word-motion-input}
 K(gb,gc)=K(b,c),\qquad g\longmapsto K(b,gc)\ \text{is continuous}
\end{equation}
for all words. These are properties of the continued distributions and the
label action on each finite level. For the boundary distributions of
Theorem~\ref{thm:joint-boundary} they are verified in
Corollary~\ref{cor:word-motion-inputs}, under the hypotheses stated there.

\begin{proposition}\label{prop:word-motion}
There is a strongly continuous unitary representation $U$ of $G$ on
$\mathcal H$ with $U(g)[b]=[gb]$. It fixes $\Omega$, preserves $\mathcal D$,
and satisfies field covariance there. There is also a linear isometry
$\mathcal V:\mathcal H\to\mathcal H_E$ with $\mathcal V[b]=q(b)$ and
$\mathcal V\Omega=\Omega_E$. Let $T_E$ be a unitary representation of the
translation group on $\mathcal H_E$, and write $\tau_ab$ for the action
of the translation by $a$ on a word $b$. Suppose that for every
translation $a$ and every word $b$ there is $\varepsilon=\varepsilon(a,b)>0$
such that the heat vectors obey
\[
 h_t(\tau_ab)=T_E(a)h_t(b)\qquad(0<t<\varepsilon).
\]
Then
$\mathcal VU(a)=T_E(a)\mathcal V$.
\end{proposition}
\begin{proof}
Extend the action on words linearly. Gram invariance makes it an isometry
of the quotient; the inverse action makes its extension to the completion
unitary. Multiplication and identity laws hold on words, hence on the dense
span, hence on the completion. For a word $b$, Gram invariance and
\eqref{eq:word-motion-input} give
\[
 \|[gb]-[g_0b]\|^2
 =2K(b,b)-2\operatorname{Re}K(g_0b,gb)\longrightarrow0.
\]
Thus orbit maps are continuous first on finite word sums and then on every
vector: approximate that vector by a finite sum, and use the isometry at both
$g$ and $g_0$ to bound the approximation errors uniformly. The vacuum and
domain statements follow on words. For a letter $a$ and a word $b$ the
equality $g(ab)=(ga)(gb)$ gives covariance on their span.

The map $\sum c_i[b_i]\mapsto\sum c_iq(b_i)$ preserves squared norms by
Lemma~\ref{lem:cold-vectors}, so is well-defined and extends isometrically.
Its vacuum identity is already known. Taking the norm limit in the stated
translation identity for heat vectors proves intertwining on every word;
boundedness then extends it to $\mathcal H$.
\end{proof}

These arguments construct all fields on one common invariant domain. They
still require the continued distributions, their locality and label-detection
properties, the continuous group action, and the independent-height limit.
The forward spectral condition and uniqueness of the invariant vacuum require
additional facts about the Euclidean translation representation and are not
consequences of the formal word construction alone.
